\paragraph{Public signals in operational forecasting.}
\citet{cui2018operational} established, in retail, that publicly
available signals (social media) measurably improve daily operational
forecasts; \citet{steinker2017value} did the same for weather in
e-commerce fulfilment. An earlier single-city study on NYC 311 by the
present author brought this template to municipal reported demand in one
city. The present work differs in scope and in kind: it spans four
heterogeneous 311 systems under one leakage-controlled protocol; adds
local-versus-pooled and temporally censored cross-city transfer; couples
forecasting to a controlled request-equivalent allocation simulation
rather than stopping at accuracy; characterizes point-forecast
non-identification, its tie-break dependence, and the family-level
served-fraction tradeoffs it exposes; and ships a stronger provenance- and
guard-backed reproducible benchmark.

\paragraph{311 analytics.}
311 data underpin a largely single-city, prediction-centric literature:
\citet{xu2017predicting} forecast Chicago 311 demand in space and time
with an adaptive kernel and motivate the exercise by resource-allocation
relevance without evaluating an allocation; \citet{cheng2022machine}
predict request volumes and their correlates in Miami-Dade County; \citet{kontokosta2017equity} and
\citet{kontokosta2021bias} document socio-spatial disparities in the
propensity to report, establishing that 311 volumes measure
\emph{reported} demand rather than need. We inherit the latter caveat
explicitly. To our knowledge no prior work evaluates 311 forecasting
across multiple cities under one protocol, nor couples it to an explicit
allocation loss with frozen pre-test environments.

\paragraph{From prediction to decisions.}
That predictive accuracy and decision quality are distinct objectives is
the founding observation of prescriptive analytics
\citep{bertsimas2020predictive} and decision-focused learning: the SPO
framework trains models against the downstream objective
\citep{elmachtoub2022smart}, task-based end-to-end learning does so in
stochastic programs \citep{donti2017task}, and \citet{mandi2024decision}
survey and benchmark the field. Those benchmarks are dominated by
synthetic costs or energy/inventory/routing settings. We contribute
complementary public-sector evidence: rather than proposing a training
method, we measure empirically where accuracy-optimized models,
distributional information, and decision-aware validation selection do
and do not change simulated allocation outcomes.

\paragraph{Capacity planning under time-varying demand.}
Setting service capacity against predictable daily variation is a
classical operations problem \citep{green2007coping}; our simulation is a
deliberately transparent discrete-time instance with integer capacity
units, family-level objectives, and carryover of unresolved
request-equivalents, designed for auditability rather than queueing
fidelity.

\paragraph{Cross-city learning and urban computing.}
Urban computing treats the city as a sensing-and-decision platform
\citep{zheng2014urban}; cross-city transfer has been studied mainly for
traffic and crowd-flow prediction \citep{wang2019crosscity}, and pooled
training across related series is standard in modern forecasting
\citep{salinas2020deepar}. We test pooled-versus-local training and
zero-shot transfer under explicit temporal censoring, and carry the
comparison through to simulated decision quality.
